\subsection{\texorpdfstring{第~\ref{sec:rigid-body}~节}{刚体}}\label{subsec:solutions-rigid-body}

\begin{description}[font=\normalfont,leftmargin=2em,labelwidth=1.5em]
\RGsolution{ex:acceleration-formula-for-planar-rigid-body}{acceleration-formula-for-planar-rigid-body}
不妨设 $m$ 支持在平面 $\xi^3=0$ 上, 并考虑测度
\[
  m_\varepsilon:=m+\varepsilon\delta_{e_3}+\varepsilon\delta_{-e_3}.
\]
因为这个测度不支持在平面上, 所以由左不变度量
\[
  \langle\!\langle V,W\rangle\!\rangle_\varepsilon
    :=\int_{\mathbb{R}^3}\langle V\xi,W\xi\rangle\,\mathrm{d}m_\varepsilon
\]
确定的协变加速度 $\nabla^\varepsilon_{\dot S}\dot S$ 满足
\[
  \langle\!\langle\nabla^\varepsilon_{\dot S}\dot S,V\rangle\!\rangle_\varepsilon
    =\int_{\mathbb{R}^3}\langle\ddot S\xi,V\xi\rangle\,\mathrm{d}m_\varepsilon.
\]
现在
\[
  \langle\!\langle V,W\rangle\!\rangle_\varepsilon
  =\int_{\mathbb{R}^3}\langle V\xi,W\xi\rangle\,\mathrm{d}m
    +2\varepsilon\langle Ve_3,We_3\rangle
\]
当 $\varepsilon\to0$ 时收敛到左不变度量
\[
  \langle\!\langle V,W\rangle\!\rangle=\int_{\mathbb{R}^3}\langle V\xi,W\xi\rangle\,\mathrm{d}m,
\]
故 $\nabla^\varepsilon_{\dot S}\dot S$ 收敛到由 $\langle\!\langle\cdot,\cdot\rangle\!\rangle$ 确定的协变加速度 $\nabla_{\dot S}\dot S$. 因而
\begin{align*}
  \langle\!\langle\nabla_{\dot S}\dot S,V\rangle\!\rangle
  &=\lim_{\varepsilon\to0}\langle\!\langle\nabla^\varepsilon_{\dot S}\dot S,V\rangle\!\rangle_\varepsilon
    =\lim_{\varepsilon\to0}\int_{\mathbb{R}^3}\langle\ddot S\xi,V\xi\rangle\,\mathrm{d}m_\varepsilon\\
  &=\lim_{\varepsilon\to0}\left(\int_{\mathbb{R}^3}\langle\ddot S\xi,V\xi\rangle\,\mathrm{d}m
                +2\varepsilon\langle\ddot Se_3,Ve_3\rangle\right)\\
  &=\int_{\mathbb{R}^3}\langle\ddot S\xi,V\xi\rangle\,\mathrm{d}m.
\end{align*}

\RGsolution{ex:prove-so-three-cross-product-isomorphism}{so-three-cross-product-isomorphism}
因为 $\mathfrak{so}(3)=\{A\in\mathfrak{gl}(3)\mid A+A^t=0\}$ 是 $3\times3$ 反对称矩阵的空间, 所以可定义映射 $\Omega:\mathfrak{so}(3)\to\mathbb{R}^3$, 它对给定的
\[
  A=\begin{pmatrix}0&-a_3&a_2\\a_3&0&-a_1\\-a_2&a_1&0\end{pmatrix}\in\mathfrak{so}(3)
\]
给出
\[
  \Omega(A)=(a_1,a_2,a_3).
\]
映射 $\Omega$ 显然为线性同构, 容易检验 $\mathfrak{so}(3)$ 上的 Lie 括号与外积等同, 即
\[
  \Omega([A,B])=\Omega(AB-BA)=\Omega(A)\times\Omega(B).
\]
此外, 给定 $\xi\in\mathbb{R}^3$, 有
\begin{align*}
  A\xi&=\begin{pmatrix}0&-a_3&a_2\\a_3&0&-a_1\\-a_2&a_1&0\end{pmatrix}
        \begin{pmatrix}\xi^1\\\xi^2\\\xi^3\end{pmatrix}
       =\begin{pmatrix}a_2\xi^3-a_3\xi^2\\a_3\xi^1-a_1\xi^3\\a_1\xi^2-a_2\xi^1\end{pmatrix}\\
     &=(a_1,a_2,a_3)\times\xi=\Omega(A)\times\xi.
\end{align*}

\RGsolution{ex:rigid-body-symmetries}{rigid-body-symmetries}
\begin{enumerate}[label=(\alph*),ref=\alph*,leftmargin=*]
\item 有
\begin{align*}
  (SIS^t)v&=S(I(S^tv))=S\int_{\mathbb{R}^3}[\xi\times(S^tv\times\xi)]\,\mathrm{d}m\\
          &=\int_{\mathbb{R}^3}S[\xi\times(S^tv\times\xi)]\,\mathrm{d}m\\
          &=\int_{\mathbb{R}^3}[(S\xi)\times(v\times(S\xi))]\,\mathrm{d}m\\
          &=\int_{\mathbb{R}^3}[\xi\times(v\times\xi)]\,\mathrm{d}m=Iv,
\end{align*}
这里利用了 $S$ 保持质量分布来作变量代换.
\item 设 $v$ 是垂直于反射平面的非零向量. 则 $Sv=-v$. 此外,
\[
  S(Iv)=SI(S^tS)v=(SIS^t)Sv=ISv=-Iv,
\]
这说明 $Iv$ 与 $v$ 平行, 因而存在一条方向为 $v$ 的主轴.
\item 设 $v$ 是一个方向与旋转轴相同的向量. 则 $Sv=S^tv=v$. 此外,
\[
  S^t(Iv)=S^t(SIS^t)v=IS^tv=Iv,
\]
这说明 $Iv$ 也被 $S^t$ 固定, 从而被 $S$ 固定. 所以 $Iv$ 与 $v$ 方向相同, 旋转轴是一条主轴.
\item 仍设 $v$ 是方向与旋转轴相同的向量. 若 $w$ 垂直于 $v$, 则 $S^tw$ 也垂直于 $v$. 此外, 若 $w$ 是 $I$ 的对应于特征值 $\alpha$ 的特征向量 (即 $Iw=\alpha w$), 则
\[
  I(S^tw)=(S^tS)IS^tw=S^t(SIS^t)w=S^tIw=\alpha S^tw,
\]
所以 $S^tw$ 也是 $I$ 的对应于 $\alpha$ 的特征向量. 考虑 $\mathbb{R}^3$ 的一组由主轴构成的基 $\{e_1,e_2,e_3\}$, 使 $e_1$ 与 $v$ 平行. 若 $S$ 不是转角为 $\pi$ 的旋转, 则 $S^te_2$ 也是特征值 $I_2$ 的特征向量, 它与 $e_2$ 线性无关且垂直于 $e_1$. 所以 $I_2$ 的特征空间包含 $\operatorname{span}\{e_2,e_3\}$. 从而每个垂直于 $v$ 的向量都是 $I_2$ 的特征向量, 所以垂直于旋转轴的所有轴都是主轴.
\end{enumerate}

\RGsolution{ex:derive-axisymmetric-kinetic-energy}{axisymmetric-kinetic-energy}
因为动能不依赖于 $\varphi$ 或 $\psi$, 不妨设 $\varphi=\psi=0$. 这样选取时, $SO(3)$ 上一条曲线的速度用 Euler 角表示为
\begin{align*}
  \dot S
  &=\dot\varphi\begin{pmatrix}0&-1&0\\1&0&0\\0&0&0\end{pmatrix}
     \begin{pmatrix}1&0&0\\0&\cos\theta&-\sin\theta\\0&\sin\theta&\cos\theta\end{pmatrix}
     +\dot\theta\begin{pmatrix}0&0&0\\0&-\sin\theta&-\cos\theta\\0&\cos\theta&-\sin\theta\end{pmatrix}\\
  &\quad+\dot\psi\begin{pmatrix}1&0&0\\0&\cos\theta&-\sin\theta\\0&\sin\theta&\cos\theta\end{pmatrix}
     \begin{pmatrix}0&-1&0\\1&0&0\\0&0&0\end{pmatrix},
\end{align*}
因而
\[
  A=S^{-1}\dot S=\dot\varphi\begin{pmatrix}0&-\cos\theta&\sin\theta\\
                   \cos\theta&0&0\\-\sin\theta&0&0\end{pmatrix}
   +\dot\theta\begin{pmatrix}0&0&0\\0&0&-1\\0&1&0\end{pmatrix}
   +\dot\psi\begin{pmatrix}0&-1&0\\1&0&0\\0&0&0\end{pmatrix},
\]
所对应的为
\begin{align*}
  \Omega&=\dot\varphi(\cos\theta\,e_3+\sin\theta\,e_2)+\dot\theta e_1+\dot\psi e_3\\
  &=\dot\theta e_1+\dot\varphi\sin\theta\,e_2+(\dot\varphi\cos\theta+\dot\psi)e_3.
\end{align*}
所以动能为
\[
  K=\frac12\langle I\Omega,\Omega\rangle
   =\frac{I_1}{2}(\dot\theta^2+\dot\varphi^2\sin^2\theta)
     +\frac{I_3}{2}(\dot\varphi\cos\theta+\dot\psi)^2.
\]

\RGsolution{ex:lagrange-top}{lagrange-top}
\begin{enumerate}[label=(\alph*),ref=\alph*,leftmargin=*]
\item 利用 $U=Mgl\cos\theta$, 容易得到运动方程:
\[
  \begin{cases}
    \dfrac{\mathrm{d}}{\mathrm{d}t}(I_1\dot\theta)
     -I_1\sin\theta\cos\theta\,\dot\varphi^2
     +I_3\sin\theta\,\dot\varphi(\dot\varphi\cos\theta+\dot\psi)=Mgl\sin\theta,\\[4pt]
    \dfrac{\mathrm{d}}{\mathrm{d}t}\bigl(I_1\dot\varphi\sin^2\theta
      +I_3\cos\theta(\dot\varphi\cos\theta+\dot\psi)\bigr)=0,\\[4pt]
    \dfrac{\mathrm{d}}{\mathrm{d}t}\bigl(I_3(\dot\varphi\cos\theta+\dot\psi)\bigr)=0.
  \end{cases}
\]
在方程中令 $\dot\theta\equiv\dot\varphi\equiv\dot\psi\equiv0$ 得到平衡点, 其条件为 $\sin\theta=0$ (假定 $l\ne0$). 这些平衡点对应于陀螺静止地处于竖直位置 (也可能倒立).
\item 容易看出, 可令 $\theta$, $\dot\varphi$ 和 $\dot\psi$ 为常数来解运动方程, 只要这些常数满足
\[
  -I_1\sin\theta\cos\theta\,\dot\varphi^2
    +I_3\sin\theta\,\dot\varphi(\dot\varphi\cos\theta+\dot\psi)=Mgl\sin\theta.
\]
若 $|\dot\varphi|\ll|\dot\psi|$, 则这些条件可近似写为
\[
  I_3\sin\theta\,\dot\varphi\dot\psi\simeq Mgl\sin\theta
   \iff\dot\varphi\simeq\frac{Mgl}{I_3\dot\psi}.
\]
\end{enumerate}
\end{description}
